\begin{originalresult}{Lemma 5.3 (minimal weight lemma; author source p.25)}
For $f:\bits^n\to\bits$ and $g:\bits^m\to\bits$, suppose that $f\comp g$ has parity certificate complexity at most $k$ at some point $x$. Let $K$ be a parameter. If the certificate complexity of $f$ at $g^n(x)$ is greater than $k\cdot K$ then we can find a blockwise partial assignment $\sigma$ to the inputs of $f\comp g$ consistent with $x$, such that for some $k'\le k$:
\begin{itemize}
\item $(f\comp g)|_\sigma$ has a parity certificate of size $k'$ at $x|_\sigma$ whose minimum weight is at least $K$,
\item $C(f|_{g^n(\sigma)},g^n(x)|_{g^n(\sigma)})>k'K$. (Note that $\sigma$ is blockwise and agrees with $x$, so here we consider the certificate complexity of $f$ at $g^n(x)$ on all the blocks that are not set by $\sigma$).
\end{itemize}
\end{originalresult}
\begin{proof}
The proof is by induction on $k$. If $k=0$, then the premises contradict each other, so we can assume that $k>0$.
Let $C$ be the parity certificate for $f\comp g$ at input $x$. By the assumption, the size of $C$ is at most $k$. If its minimum weight is at least $K$, then we are done. Otherwise, without loss of generality, $C$ contains a nontrivial constraint $\ell=\alpha$ whose support is less than $K$. Choose a partial assignment that agrees with $x$ on the support of this constraint, and complete it to a blockwise partial assignment $\sigma$ consistent with $x$ (complete only the blocks with at least one variable assigned). Also, let $\tau:=g^n(\sigma)$.
By construction, $(f\comp g)|_\sigma$ has a parity certificate of size at most $k-1$ at $x|_\sigma$, namely $C\setminus\{\ell=\alpha\}$. On the other hand, since we only set the values of at most $K$ blocks, the certificate complexity of $f$ reduces by at most $K$ after being restricted to $\tau$, and hence $C(f|_\tau,g^n(x)|_\tau)>(k-1)K$. Now we can apply the induction hypothesis for $k-1$.
\end{proof}
The following lemma demonstrates an essential property of the parity certificates with huge minimal weight:
\begin{originalresult}{Lemma 5.4 (author source pp.25--26)}
Suppose that $C_\oplus$ is a parity certificate with minimum weight at least $K$ and $C$ is a (regular) certificate of size less than $K$. Then there is a point that satisfies both $C$ and $C_\oplus$.
\end{originalresult}
In the proof, the parity certificate is processed with a careful Gaussian elimination, and then the variables fixed by the certificate $C$ are free variables of the parity certificate.
\begin{proof}
Let $V+b$ be the affine subspace corresponding to $C_\oplus$, where $V$ is a linear subspace. We pick a basis $\ell_1,\ldots,\ell_k$ for $V$ and indices $i_1,\ldots,i_k\in[n]$ such that:
\begin{itemize}
\item $z_{i_t}$ is in the support of $\ell_t$,
\item $z_{i_s}$ is not in the support of $\ell_t$ for $s<t$,
\item $z_{i_t}$ does not belong to the support of $C$.
\end{itemize}
We do it inductively. At step $t$, we have already picked $\ell_1,\ldots,\ell_{t-1}$. Consider vectors in $V$ whose support does not contain any of the variables $z_{i_1},\ldots,z_{i_{t-1}}$. Since $V$ has minimum weight $K$, we can find such a vector $\ell_t$ whose support contains a variable $z_{i_t}$ with $z_{i_t}\notin C$. By the assumption, this vector does not contain any of the variables $z_{i_1},\ldots,z_{i_{t-1}}$.
Now we can construct an input $\alpha$ which satisfies both $C$ and $C_\oplus$. First, we set the variables from $C$ such that the resulting input satisfies $C$. Then we set arbitrarily all the variables not occurring in $C$ or $\{z_{i_1},\ldots,z_{i_k}\}$. Finally, we set variables $z_{i_k},\ldots,z_{i_1}$ (in this order) in such a way that the resulting input satisfies $C_\oplus$.
\end{proof}
